\chapter{Roses and fabric plots}\label{ch:roses}
\begin{goals}
\textbf{In this chapter} you draw rose diagrams of strikes, dip directions and trends, and fabric plots
that place whole datasets on Vollmer's or Woodcock's diagram, and you choose how each is binned and drawn.
\end{goals}

A net shows every measurement. Two other kinds of diagram summarize: a \term{rose diagram} counts azimuths
in sectors, and a \term{fabric plot} reduces each dataset to one point by the shape of its fabric. They are
diagrams like a net: each has its own tab, its own layers and its own settings, and they show the same
datasets as everything else (\cref{sec:data-once}).

\section{Making one}

Choose \menu{Diagram>New rose diagram} or \menu{Diagram>New fabric plot}, or type \code{new rose} or
\code{new fabric} in the console. The new diagram starts with views of every layer on the active diagram,
and opens in a tab of its own: \ui{Rose 1}, \ui{Fabric 1}. To add more data later, drag a dataset or a
layer onto it in the tree, as for a net (\cref{ch:diagrams}).

A rose or a fabric plot draws only its \emph{visible} layers, so tick and untick layers in the tree to
change what it counts. The settings are in the diagram's own menu: right-click its name in the tree.
Rotation, projection and hemisphere belong to nets and don't apply here.

\section{Rose diagrams}\label{sec:rose}

\shot[0.6\linewidth]{rose-strike}{A rose of the Tocher sample's strikes in 10° sectors, area-true. Strike is
an axis, so each petal has its twin opposite.}

A rose bins the azimuths of every visible layer together, in one colour (the first layer's), with four
reference rings. What it bins is the diagram's \ui{Quantity}:

\begin{description}
\item[\ui{strike (planes, axial)}] the strike of each plane, by the right-hand rule (dip direction less
90°). A strike of 030 and one of 210 are the same line on the map, so the rose counts each strike at both
ends: the petals come in opposite pairs. This is the default.
\item[\ui{dip direction (planes)}] the dip direction of each plane.
\item[\ui{trend (lines)}] the trend of each line.
\end{description}

Planes have no trend and lines no strike, so a quantity skips the measurements it doesn't apply to. If
none are left, the diagram says so instead of drawing: \ui{no lines to bin by trend — choose another
quantity in the diagram's menu}.

A dip direction or a trend points one way, so their roses count each azimuth once: a set dipping to 120
has a petal at 120 and none at 300. When you want them read as axes anyway (lineations whose sense means
nothing, say), tick \ui{Count both ends (axial)} in the diagram's menu. A strike rose is always axial,
and the item is greyed out there.

Two more settings, also in the diagram's menu:

\begin{description}
\item[\ui{Sector}] the width of each sector: 5°, 10° (the default), 15°, 20° or 30°.
\item[\ui{Scale}] \ui{area-true (√)}, the default, makes a petal's \emph{area} proportional to its count,
so the radius grows as the square root. \ui{linear} makes the radius proportional to the count, which
exaggerates the larger petals.
\end{description}

\begin{why}
A petal's area is what the eye compares. On a linear rose a sector with four times the count has sixteen
times the area, and the largest petal dominates the picture far beyond its share. Area-true scaling keeps
what you see in proportion to what was measured.
\end{why}

\section{Fabric plots}\label{sec:fabric}

\shot[0.6\linewidth]{fabric-vollmer}{A Vollmer plot of the Tocher sample and the two sets it splits into by
dip direction (\cref{sec:stats-intersect}): \code{west} is a cluster, \code{east} and the whole a girdle.}

A fabric plot places each visible layer with at least three measurements as one point, in its own colour
and labelled with its name. The point comes from the layer's eigenvalues, the same ones the statistics card
shows (\cref{sec:stats-eigen}). The diagram's \ui{Plot} setting chooses the diagram:

\begin{description}
\item[\ui{Vollmer P–G–R}] (the default) a triangle with corners for a perfect point (cluster), girdle and
random fabric \parencite{vollmer1990}. A dataset sits at its P, G and R fractions: Tocher's 0.16, 0.63,
0.21 puts it near the girdle corner.
\item[\ui{Woodcock K–C}] $\ln(S_2/S_3)$ across and $\ln(S_1/S_2)$ up \parencite{woodcock1977}. Clusters
fall above the line $K = 1$ and girdles below it; the further from the origin, the stronger the fabric.
\end{description}

The fabric plot is drawn in dark ink on a light card in either theme, as it would be on paper.

For the three sets of \cref{fig:fabric-vollmer}:

\begin{center}\small
\begin{tabularx}{\linewidth}{@{}lrrrrL@{}}\toprule
\sffamily\bfseries set & \sffamily\bfseries P & \sffamily\bfseries G & \sffamily\bfseries R & \sffamily\bfseries K & \sffamily\bfseries reads as \\\midrule
\code{tocher} & 0.16 & 0.63 & 0.21 & 0.21 & a girdle \\
\code{east} & 0.15 & 0.75 & 0.11 & 0.12 & a stronger girdle \\
\code{west} & 0.64 & 0.21 & 0.15 & 1.46 & a cluster \\
\bottomrule\end{tabularx}
\end{center}

\begin{tip}
A fabric plot is the quickest way to compare domains. Split a dataset into structural domains with
selections (\cref{sec:select-dataset}), draw them all on one fabric plot, and the ones that behave
differently stand apart.
\end{tip}

\section{In the files}

A rose or a fabric plot is a diagram like a net, so \file{/project} has it as a directory
(\cref{ch:project-fs}): \file{diagrams/Rose 1/diagram.json} holds its settings
(\code{"quantity": "strike"}, \code{"bins": 10}, \code{"scale": "area"}, or \code{"plot": "vollmer"}), and
\file{diagrams/Rose 1/rose.svg} or \file{diagrams/Fabric 1/fabric.svg} is the diagram as drawn, ready to
copy out. Writing \file{diagram.json} changes the settings, as the menu does.

\begin{tryit}
\begin{enumerate}
\item Open the Tocher sample and choose \menu{Diagram>New rose diagram}. A rose of strikes opens in a new
tab.
\item In the tree, right-click \ui{Rose 1} and choose \ui{Sector › 30°}. The petals widen; choose
\ui{Scale › linear} and see the largest ones grow.
\item In the console, type \code{new fabric}. A Vollmer triangle opens with one point, \code{tocher},
near the girdle corner.
\item Type \code{ls "diagrams/Fabric 1"} and \code{head -3 "diagrams/Fabric 1/fabric.svg"}: the plot as a
file.
\end{enumerate}
\end{tryit}
